\chapter*{Záver}  % chapter* je necislovana kapitola
\addcontentsline{toc}{chapter}{Záver} % rucne pridanie do obsahu
\markboth{Záver}{Záver} % vyriesenie hlaviciek

Výsledky tejto práce ukazujú, že navrhnutý autonómny agent pre hru Trackmania dosahuje konkurenčný výkon v porovnaní s existujúcimi riešeniami ako Trackmania AI a TMRL. Na základe experimentov sa ukázalo, že agent dokáže prejsť rôzne trate s vysokou konzistenciou a rýchlosťou odozvy na dynamické zmeny v prostredí. V porovnaní s konkurenčnými projektmi, náš agent dosiahol porovnateľné časy prejdenia trate, pričom vykazoval vyššiu spoľahlivosť a rýchlejšiu odozvu.

Hlavné prednosti nášho agenta zahŕňajú plnú autonómnosť bez potreby ľudskej interakcie, rýchlu odozvu a flexibilitu do budúcnosti, ktorá umožňuje rozšírenie jeho schopností na trate s výškovým prevýšením alebo slučkami. Tieto výsledky potvrdzujú efektívnosť zvoleného prístupu a poskytujú pevný základ pre ďalší výskum a zlepšovanie autonómneho riadenia v prostredí hry Trackmania.

Napriek dosiahnutým úspechom existujú oblasti, ktoré by si zaslúžili ďalšiu pozornosť. Bolo by vhodné preskúmať mapy s alternatívnymi cestami, trénovať agenta na rôznych povrchoch a naučiť ho jazdiť aj na tratiach s výškovými rozdielmi. Ďalším smerom by mohla byť optimalizácia jazdy agenta na dosahovanie nadľudských výkonov, čo zahŕňa zlepšenie stratégie riadenia, zníženie chýb a zvýšenie konzistencie výkonov na rôznych typoch tratí. 

Táto práca ukazuje, že aj v herných prostrediach je možné dosiahnuť vysokú úroveň autonómie a výkonnosti pomocou moderných techník strojového učenia. Dúfame, že naše výsledky inšpirujú ďalší vývoj v oblasti autonómnych agentov v hre Trackmania.

